% =====================================================================
%  夏ゼミ配布資料（レジュメ） 松本 鋭
%  研究室配布の template.tex（jarticle・2段組）をそのまま土台にしている。
%  コンパイル: platex natsuzemi_resume_matsumoto.tex （表と参考文献の
%              番号を確定させるため2回）→ dvipdfmx natsuzemi_resume_matsumoto.dvi
%  ※ 日付は開催日に置き換える
% =====================================================================
\documentclass[a4j,twocolumn]{jarticle}
\usepackage{fancybox,wrapfig,ascmac,tabularx}
\usepackage[dvipdfmx]{graphicx}

\makeatletter
\def\section{\@startsection{section}{1}{\z@}{2ex plus .2ex minus .2ex}%
			{.5ex plus .2ex minus .2ex}{\large\bfseries}}
\def\thesection{\arabic{section}.}
\def\subsection{\@startsection{subsection}{1}{\z@}{.7ex plus .2ex minus .2ex}%
			{.5ex plus .2ex minus .2ex}{\normalsize\bfseries}}
\def\thesubsection{\arabic{section}.\arabic{subsection}}
\def\thefootnote{\fnsymbol{footnote}}
\makeatother


\def\baselinestretch{0.8}

\setlength{\textheight}{23.5cm}%297-30-27 - 5
\setlength{\textwidth}{17.4cm}%210-18-18 - 10
%\setlength{\headheight}{0.0in}
\setlength{\headsep}{0.0in}
\setlength{\oddsidemargin}{-.9cm}%+3
\setlength{\evensidemargin}{-.9cm}%+3
\setlength{\columnsep}{7mm}

% local settings
\usepackage{amsmath,amssymb}
\usepackage{bm}
% end of local settings

\begin{document}
\pagestyle{empty}
\thispagestyle{empty}

\twocolumn[%

\begin{flushright}
2026/08/XX 夏ゼミ配布資料
\end{flushright}

\begin{center}
 {\Large SELDモデルを用いた\\
  難聴歩行者向け屋外危険音通知システムの構築}\vspace{.5ex}

\large
\mbox{}
\hfil
\setcounter{footnote}{1}
{\bfseries 松本 鋭*}
\hfil
\mbox{}

\mbox{}
\hfil
{\sffamily　Satoshi Matsumoto}
\hfil
\mbox{}
\hfil

\end{center}
]

\setcounter{footnote}{1}
\footnotetext{東京理科大学 創域理工学部 情報計算科学科 4年}

% ---------------------------------------------------------------
\section{背景と目的}

\subsection{背景}
歩行者は，接近する車やサイレン，踏切といった危険の察知を聴覚に大きく
依存している．難聴者はこの情報が欠落するため，屋外歩行における事故リスクが
高い．当事者調査\cite{findlater}でも，緊急車両などの危険音は「通知して
ほしい音」の上位に挙げられ，なかでも\textbf{音がどちらから来るか}が
最も重視される情報だと報告されている．同じ課題は，イヤホンを装着した
歩行者や高齢者にも共通する．

一方で「機械の耳」は疲労せず全方位を常時監視でき，センサが対応すれば
人間の可聴域外の信号も扱える．したがって，失われた聴覚をそのまま補うので
はなく，人間の聴覚が届かない範囲まで含めた支援が期待できる．

\subsection{目的}
本研究の目的は，屋外の危険音8クラス（サイレン・クラクション・バック音・
自転車ベル・車（EV・大型含む）・踏切と列車・キックボード・バイク）に
ついて\textbf{種類・方向・距離}を同時に推定し，それを\textbf{危険度3段階の
通知}に変換して届けるシステムを構築することである．

距離を含めるのは，鳴らすかどうかの判断が距離で決まるためである．
最終的には首元の振動で伝える構成を想定するが，システムは「何がどこに
いるか」を伝えるにとどめ，行動の判断は利用者に委ねる．

% ---------------------------------------------------------------
\section{基礎知識}

\subsection{SELD}
SELD（Sound Event Localization and Detection）は，画像における物体検出の
音・方向版にあたるタスクであり，何がいつ鳴っているかを当てる音響イベント
検出（SED）と，どの方向から鳴っているかを当てる音源定位（SSL）を，複数の
音源について同時に解く\cite{adavanne}．

\subsection{音源距離推定}
本研究では SELD に音源距離推定（SDE）を加え，種類・方向・距離の同時推定へ
拡張する．距離は通知の危険度を決める量であり，「自分に近づいてくる車」と
「遠方を通り過ぎる車」を区別するために不可欠である．

% ---------------------------------------------------------------
\section{関連研究}

\subsection{PSELDNets（用いる基盤モデル）}
Hu ら\cite{pseldnets}は，1,167時間・170クラスの大規模合成データで事前学習
した SELD の基盤モデル PSELDNets を提案し，主要データセットで最高性能を
達成した．FOA形式の4ch空間音響を入力とし，同一クラスの音源を最大3つまで
同時に検出できる multi-ACCDOA\cite{maccdoa} 出力を持つ．本研究は，これに
距離推定のヘッドを加えて微調整する土台として用いる．

\subsection{DynamicSound（用いる屋外シミュレータ）}
Barbisan ら\cite{dynamicsound}は，室内残響ツールでは扱えなかった屋外・
移動音源の状況を，伝搬遅延・ドップラー・距離減衰・大気吸収・地面反射と
いう物理で明示的にモデル化した屋外音響シミュレータを公開している．ただし
提供されるのは物理的な音の生成までで，SELD の学習に必要な FOA 化・ラベル
付け・評価は含まれない．調査した範囲では，これを SELD の学習に用いた公開例
も見当たらない．

\subsection{本研究の位置づけ}
\textbf{歩行者装着・複数の危険クラス同時・距離まで}を屋外で扱った研究は，
調査した範囲では見当たらない．また，音の分類と方向の推定を別々の手法に
分業させる構成は，複数の音源が同時に鳴ったときに「どの音がどの方向か」の
対応付けが壊れやすい．そこで本研究では，両者を同時に解く SELD を採用する．

% ---------------------------------------------------------------
\section{提案手法}

\subsection{合成データの自作}
屋外 SELD の最大の障害は，本研究の条件（屋外・歩行者視点・距離ラベル
付き）を満たす学習データが存在しないことである．そこで物理シミュレー
ションを自作し，4ch の音と方向・距離のラベルを同一の計算経路から同時に
出力する．音源の音量・周波数は日本の法規・省令・実測研究に出典を持たせた．

\subsection{2層構造}
知覚層は8クラスの検出・方向・距離を0.1秒ごとに推定し，通知層はその推定から
「いつ・どの強さで」伝えるかを決める．危険のたびに全て鳴らすと実用に耐え
ないため，\textbf{通知の頻度そのものを設計対象}としている点が特徴である．

\subsection{通知層の改良：最接近距離の予測}
当初の通知層は推定距離のしきい値のみで判定しており，速度を見ていないため
相手が速いほど余裕が縮む弱点を持つ．そこで，距離の縮み方 $\dot d$ と方位の
振れ方 $\dot a$ から\textbf{最接近の距離と時刻}を等速直線運動として解く方式
に改めた．相対速度の大きさが $|\bm v|^{2}=\dot d^{2}+(d\dot a)^{2}$ と
書けることから，到達時刻 $t_{\mathrm{cpa}}$ と予測最接近距離
$d_{\mathrm{cpa}}$ は
\begin{equation}
  t_{\mathrm{cpa}}=\frac{-d\,\dot d}{|\bm v|^{2}},\qquad
  d_{\mathrm{cpa}}=\frac{d^{2}|\dot a|}{|\bm v|}
\end{equation}
となる．$d_{\mathrm{cpa}}$ が小さく $t_{\mathrm{cpa}}$ も短いときに警告
する．これにより正面から接近する対象には遠方で鳴り，横を通り過ぎるだけの
対象には鳴らない．

% ---------------------------------------------------------------
\section{実験内容と結果}

\subsection{実験内容}
知覚層は PSELDNets を距離ヘッドごと微調整し，自作の合成10,200クリップ
（8クラス・約28時間）で学習した．評価は，学習にも検証にも使っていない
新規1,800クリップを，基準を先に決めて\textbf{1回だけ}採点する事前登録の
手順による．指標は SELD の標準的な検出率・F値・方向誤差\cite{mesaros}に，
距離の絶対誤差と通知の到達率を加えた．

\begin{table}[t]
\caption{評価用1,800クリップでの結果．通知層は1.5m以内まで接近する車
508台が分母で，従来＝距離しきい値のみ，提案＝最接近距離の予測である}
\label{tab:result}
\centering\footnotesize
\begin{tabular}{lrr}
\hline
\multicolumn{3}{l}{\textbf{知覚層}} \\
検出率（8クラス） & \multicolumn{2}{r}{97.2--100\%} \\
総合F値（位置込み） & \multicolumn{2}{r}{91.3\%} \\
方向誤差の中央値 & \multicolumn{2}{r}{1--5度} \\
距離の絶対誤差（至近5m以内） & \multicolumn{2}{r}{0.21m} \\
静音EVのみの場面の検出率 & \multicolumn{2}{r}{98.99\%} \\
\hline
\multicolumn{1}{l}{\textbf{通知層}} & 従来 & 提案 \\
至近警告が届いた割合 & 69.9\% & 85.0\% \\
警告から最接近までの余裕 & 0.00秒 & 0.80秒 \\
安全な車に鳴らさなかった割合 & 90.4\% & 85.2\% \\
\hline
\end{tabular}
\end{table}

\subsection{実験結果}
結果を表\ref{tab:result}に示す．方向誤差が幅を持つのは，出力・正解ともに
1度刻みで記録しておりこれ以上分解できないためである．距離は相対誤差でみると
1.5--30m の範囲でおおむね 6--9\% と安定している．

% ---------------------------------------------------------------
\section{評価と考察}

最接近距離の予測に変えたことで，至近警告が届いた割合は 15.1 ポイント
向上した．とくに重要なのは余裕時間で，従来方式の中央値 0.00 秒は
\textbf{警告が最接近と同時，すなわち間に合っていなかった}ことを意味する．
一方で安全な車に鳴らさずに済んだ割合は 5.2 ポイント低下した．早く鳴らすことと
余計に鳴らさないことは両立せず，この対価は正直に報告すべきである．

比較のため，接近速度のみを用いる到達時間（TTC）\cite{vogel}による判定も
試したが，横を通り過ぎるだけの対象にも警告が出るため，安全な対象を鳴らさ
ずに済んだ割合が 82.1\% から 15.7\% へ崩壊した．接近しているかだけでは
不十分であり，\textbf{進行方向の情報が不可欠}であることが確認できた．

限界として，以上はすべて合成データ上の評価であり，車の至近帯は安全上再現
できないため，実環境での確認は自転車とキックボードによる部分検証にとどまる．

% ---------------------------------------------------------------
\section{今後の方針}

\begin{itemize}\itemsep=0pt \parskip=0pt
  \item \textbf{実環境録音による検証}：4ch FOA マイクを胸部に装着し，統制
        録音・機会捕捉・負例の3層で100本，歩行の影響を測る静止／歩行の対で
        100本を収録する．
  \item \textbf{物理 ablation}：3.2 節の物理を1つずつ外した4条件を同一の
        土俵で採点し，「どの物理が実環境性能に必要か」に直接答える．
  \item \textbf{距離推定の精度向上}：上限実験から，余裕時間を伸ばす鍵が
        距離推定の精度と複数音源の追跡にあると分かっている．
\end{itemize}

% ---------------------------------------------------------------
\begin{thebibliography}{99}\footnotesize
\setlength{\itemsep}{0pt}
\bibitem{pseldnets} J.~Hu \textit{et al.}, ``PSELDNets: Pre-trained Neural
  Networks on a Large-scale Synthetic Dataset for Sound Event Localization
  and Detection,'' IEEE/ACM Trans. Audio, Speech, Language Process.,
  arXiv:2411.06399, 2024.
\bibitem{dynamicsound} L.~Barbisan, M.~Levorato, F.~Riente, ``DynamicSound
  simulator for simulating moving sources and microphone arrays,''
  arXiv:2601.15433, 2026.
\bibitem{maccdoa} K.~Shimada \textit{et al.}, ``Multi-ACCDOA: Localizing and
  Detecting Overlapping Sounds from the Same Class with Auxiliary
  Duplicating Permutation Invariant Training,'' Proc. ICASSP, 2022.
\bibitem{adavanne} S.~Adavanne \textit{et al.}, ``Sound Event Localization
  and Detection of Overlapping Sources Using Convolutional Recurrent Neural
  Networks,'' IEEE J. Sel. Top. Signal Process., vol.~13, no.~1,
  pp.~34--48, 2019.
\bibitem{mesaros} A.~Mesaros \textit{et al.}, ``Joint Measurement of
  Localization and Detection of Sound Events,'' Proc. WASPAA,
  pp.~333--337, 2019.
\bibitem{findlater} L.~Findlater \textit{et al.}, ``Deaf and Hard-of-hearing
  Individuals' Preferences for Wearable and Mobile Sound Awareness
  Technologies,'' Proc. CHI, 2019.
\bibitem{vogel} K.~Vogel, ``A comparison of headway and time to collision as
  safety indicators,'' Accident Analysis \& Prevention, vol.~35, no.~3,
  pp.~427--433, 2003.
\end{thebibliography}

\end{document}
